% ===================================================================
%  TABELUL Nr. 2 — Corpul omenesc. — Recreația. — Jocurile.
%  Traduction 2026 d'apres texto/io, controlee sur texto/fr, texto/en
%  et texto/es. Consigne : voir texto/ro/10-tabelul-01.tex.
%
%  LE TABLEAU DE L'ANATOMIE MONTRE LE PARTAGE ANNONCE AU TABLEAU 1 :
%  le corps est latin presque partout — « cap », « braț », « picior »,
%  « inima », « ochi », « mâna » — mais la joue est « obraz », qui est
%  slave, a cote du front « frunte », qui est latin. Deux couches dans
%  le meme visage, et c'est le cas partout dans cette colonne.
% ===================================================================

%%K t02-apar-1 apar
\VUsaut{4.0mm}
\VUtitre{15.0pt}{60}{TABELUL Nr. 2}
\VUsaut{2.0mm}
\VUtitre{11.4pt}{0}{Corpul omenesc. --- Recreația.}
\VUtitre{11.4pt}{0}{Jocurile.}

%%K t02-tit-0 sub
\VUtitre{13.2pt}{0}{Corpul omenesc.}
\VUsaut{2.4mm}
\VUcentre{10.2pt}{0}{\textit{(Lecție simplă de științe naturale.)}}
\VUsaut{3.0mm}

%%K t02-01-1 p
1. --- Părțile de căpetenie ale corpului omenesc sunt: \VUgras{capul},
\VUgras{trunchiul} și \VUgras{membrele}.

%%K t02-tit-1 sub
\VUpk{10.2pt}{40}{Capul.}
\VUsaut{3.0mm}

%%K t02-02-1 p
2. --- Capul are două părți: \VUgras{craniul}, în care se află \VUgras{creierul}
și care este acoperit de \VUgras{păr}, și \VUgras{fața}.

%%K t02-03-1 p
3. --- În desenul nostru nu se văd nici craniul, nici creierul; în schimb se văd
foarte limpede părțile însemnate ale feței.

%%K t02-04-1 p
4. --- Iată mai întâi \VUgras{fruntea}, care vorbește despre o minte puternică și
despre un gând necontenit la lucru. \VUgras{Tâmplele} sunt foarte descoperite.
\VUgras{Ochiul} este viu și larg deschis. O \VUgras{sprânceană} deasă și
\VUgras{gene} lungi îl apără.

%%K t02-05-1 p
5. --- Iată de asemenea \VUgras{nasul} și \VUgras{nările}. Nasul ne slujește
pentru miros și pentru răsuflare. De o parte și de alta a nasului sunt
\VUgras{obrajii}, iar mai departe \VUgras{urechile}. Partea de jos a feței o
alcătuiesc \VUgras{gura} și \VUgras{bărbia}.

%%K t02-06-1 p
6. --- Se văd rău, fiindcă \VUgras{mustața} și \VUgras{barba} ni le ascund. Dar
știm că în gură sunt două \VUgras{buze}, \VUgras{maxilarul de sus} și
\VUgras{maxilarul de jos} cu \VUgras{dinții} lor, și \VUgras{limba}. Cu gura
vorbim și mâncăm. Cu limba simțim gustul mâncării.

%%K t02-07-1 p
7. --- Ochiul, urechea, nasul și gura, împreună cu degetele, sunt organele celor
cinci simțuri: \VUgras{văzul}, \VUgras{auzul}, \VUgras{mirosul},
\VUgras{gustul} și \VUgras{pipăitul}.

%%K t02-08-1 p
8. --- Capul este așadar una dintre părțile cele mai curioase ale corpului
omenesc.

%%K t02-tit-2 sub
\VUsaut{4.4mm}
\VUpk{10.2pt}{40}{Trunchiul.}
\VUsaut{3.0mm}

%%K t02-09-1 p
9. --- \VUgras{Gâtul} leagă capul de trunchi. În el intră \VUgras{beregata} și
\VUgras{ceafa}.

%%K t02-10-1 p
10. --- Și în trunchi se află multe organe de neapărată trebuință. În
\VUgras{piept} se găsesc \VUgras{plămânii}, cu care răsuflăm, și \VUgras{inima},
miezul vieții. Când inima încetează să bată, ființa vie moare.

%%K t02-11-1 p
11. --- \VUgras{Coastele} țin pieptul.

%%K t02-12-1 p
12. --- Celelalte părți ale trunchiului sunt \VUgras{coasta}, \VUgras{spatele},
\VUgras{umerii} și \VUgras{pântecele} (burta). Ne culcăm pe o coastă sau pe spate,
iar greutățile le purtăm pe umăr.

%%K t02-13-1 p
13. --- În pântece se află \VUgras{stomacul} și \VUgras{mațele}. Ne slujesc la
mistuirea mâncării.

%%K t02-tit-3 sub
\VUsaut{4.4mm}
\VUpk{10.2pt}{40}{Membrele.}
\VUsaut{3.0mm}

%%K t02-14-1 p
14. --- Avem patru \VUgras{membre}: două \VUgras{brațe} și două
\VUgras{picioare}. Cu brațele luăm, ținem, ridicăm și adunăm lucrurile; cu
picioarele stăm, umblăm și alergăm.

%%K t02-15-1 p
15. --- \VUgras{Umărul} leagă brațul de corp. Un \VUgras{mușchi} foarte puternic,
cel cu două capete, mișcă brațul, pe care \VUgras{cotul} îl leagă de
\VUgras{antebraț}. \VUgras{Încheietura mâinii} alcătuiește articulația
\VUgras{mâinii}, care se sfârșește cu \VUgras{degetele} și cu \VUgras{unghiile}.
Pe mână se vede \VUgras{vena} (și artera), prin care curge \VUgras{sângele}.

%%K t02-16-1 p
16. --- Părțile piciorului sunt: \VUgras{coapsa}, \VUgras{genunchiul} și
\VUgras{scobitura genunchiului}, \VUgras{pulpa}, a cărei \VUgras{carne} este
acoperită de \VUgras{piele}, \VUgras{glezna} și \VUgras{laba piciorului}.
\VUgras{Oasele} piciorului sunt foarte groase și foarte tari. La capătul labei
sunt \VUgras{degetele de la picioare}. Celelalte părți sunt \VUgras{talpa} și
\VUgras{călcâiul}.

%%K t02-17-1 p
17. --- Aceasta este descrierea aproape întreagă a corpului omenesc.

%%K t02-17-2 p
\textit{Observație.} --- \textit{Cu ajutorul întorsăturii «mă doare... (capul și
așa mai departe)» profesorul va putea străbate din nou tot acest vocabular din
partea} \VUgras{sănătății trupești} \textit{bune sau rele}.

%%K t02-tit-4 sub
\VUsaut{5.0mm}
\VUpk{10.2pt}{40}{Gimnastica.}
\VUsaut{2.4mm}
\VUcentre{10.2pt}{0}{\textit{(Povestită de unul dintre elevi.)}}
\VUsaut{3.0mm}

%%K t02-18-1 p
18. --- Ca să fim mlădioși, îndemânatici și sănătoși, trebuie să ne deprindem
picioarele și brațele. Iată de ce ne jucăm și facem \VUgras{gimnastică}.

%%K t02-19-1 p
19. --- În timpul săptămânii amândouă aceste îndeletniciri se petrec în curtea
școlii (vezi tabelul nr. 1). Dar joia și duminica ieșim pe o pajiște mare, unde
avem loc să alergăm și să ne veselim.

%%K t02-20-1 p
20. --- Profesorii și părinții noștri vin uneori cu noi; dar exercițiile le
conduce \VUgras{profesorul de gimnastică} \textsuperscript{(12)}.

%%K t02-21-1 p
21. --- Pe pajiște, la stânga intrării, se află \VUgras{aparatele de gimnastică}
\textsuperscript{(1)}: ne cățărăm pe \VUgras{scară} \textsuperscript{(2)}, ne
legănăm cu capul în jos pe \VUgras{inele} \textsuperscript{(3)} și pe
\VUgras{trapez} \textsuperscript{(4)}, ne ridicăm în brațe până în vârful
\VUgras{scării de frânghie} \textsuperscript{(5)} sau al \VUgras{frânghiei cu
noduri} \textsuperscript{(6)}.

%%K t02-22-1 p
22. --- În vremea aceasta tovarășii noștri sar peste \VUgras{calul de lemn}
\textsuperscript{(7)} sau peste \VUgras{paralele} \textsuperscript{(11)}. Alții se
\VUgras{boxează} \textsuperscript{(8)} sau fac \VUgras{scrimă}
\textsuperscript{(9)} cu masca și cu \VUgras{floreta}. Alții, în sfârșit, ridică
\VUgras{halterele} \textsuperscript{(10)}.

%%K t02-tit-5 sub
\VUsaut{4.6mm}
\VUpk{10.2pt}{40}{Jocurile.}
\VUsaut{3.0mm}

%%K t02-23-1 p
23. --- În jurul gimnaștilor băieții și fetele se joacă de-a fel de fel de
\VUgras{jocuri}. Fetele se desfată cu \VUgras{cercul} lor \textsuperscript{(14)};
sar \VUgras{coarda} \textsuperscript{(13)} sau își cheamă frații și verișoarele la
o partidă de \VUgras{crochet} \textsuperscript{(15)}. Ce bucuroase sunt când
izbesc \VUgras{bila} potrivnicului cu \VUgras{ciocanul lor de lemn} tocmai când
acesta crede că va trece ușor prin portiță!

%%K t02-24-1 p
24. --- Băieții preferă jocurile mai vioaie. Se joacă bucuros de-a
\VUgras{popicele} \textsuperscript{(16)}, pe care le doboară cu îndemânare cu
\VUgras{bila} \textsuperscript{(17)}. Nu disprețuiesc nici \VUgras{titirezul}
\textsuperscript{(18)}, nici \VUgras{bilboquet-ul} \textsuperscript{(19)}, și nici
\VUgras{broasca} \textsuperscript{(20)}. Dar jocurile lor cele mai dragi sunt
\VUgras{bilele} \textsuperscript{(21)} și \VUgras{mingea de perete}
\textsuperscript{(22)}, fiindcă zidul \VUgras{serei} \textsuperscript{(26)} este
tocmai bun ca să întoarcă mingea ușoară.

%%K t02-25-1 p
25. --- Micul Ioan este foarte îndemânatic pe \VUgras{picioroangele} lui
\textsuperscript{(24)} și își ține minunat cumpăna. Lângă el câțiva tovarăși se
joacă de-a \VUgras{ascunselea} \textsuperscript{(25)} în seră și în spatele
\VUgras{gardului viu}. Alții preferă \VUgras{baba-oarba pe mână}
\textsuperscript{(28)} sau \VUgras{fluturașul} \textsuperscript{(29)}, care zboară
de la o \VUgras{rachetă} la alta.

%%K t02-noto-1 noto
\VUnotes{143.2mm}{%
(*) Jocurile copiilor poartă adesea nume cu totul naționale. Le-am numit în ido
cum am putut: când după însăși fapta care le deosebește, când luând numele
național al unei țări sau al alteia, atunci când nu tulbură prea mult. De
pildă: \VUgras{butoiul} (nemțesc și franțuzesc) este \VUgras{broasca}
\textsuperscript{(20)} (englezesc), unde jucătorii se silesc să arunce jetoanele
rotunde în gura unei broaște de metal, căscată deasupra unei lăzi găurite
(butoi). \VUgras{Săritura peste umeri} \textsuperscript{(30)} se deosebește de
\VUgras{capra} numai prin aceea că, sărind peste cap, jucătorii se sprijină cu
mâinile pe umerii tovarășului, pe când la «\VUgras{capră}» se sprijină numai pe
spate. --- Sau iarăși: unii jucători se apleacă, iar ceilalți sar peste spinările
lor, sprijinindu-se cu mâinile pe umeri; cei dintâi trebuie să rabde izbitura
fără să se clatine, cei de-ai doilea să sară fără să-și piardă cumpăna (vezi
tabelul însuși).%
}

%%K t02-26-1 p
26. --- În mijlocul pajiștii cresc doi copaci. La poalele unuia dintre ei șapte
sau opt băieți au început un joc de \VUgras{săritură peste umeri}
\textsuperscript{(30)} (*). Acest joc nu este cunoscut în toate țările; dar
oricine știe ce este \VUgras{capra}, de care copiii acum nu se mai joacă.

%%K t02-tit-6 sub
\VUsaut{5.0mm}
\VUpk{10.2pt}{40}{Jocuri în aer liber.}
\VUsaut{3.4mm}

%%K t02-27-1 p
27. --- \VUgras{Baba-oarba} \textsuperscript{(31)} este de asemenea un joc foarte
vesel; fetele și băieții își pun pe rând \VUgras{legătura} la ochi și îi prind pe
cei greoi, care se lasă prinși.

%%K t02-28-1 p
28. --- În mijlocul pajiștii --- iată un joc cu totul franțuzesc, deși cam
grosolan: este \VUgras{ursul} \textsuperscript{(32)}, mult mai zgomotos decât
pașnicul joc al \VUgras{zmeului} \textsuperscript{(33)} și chiar decât jocul
englezesc de \VUgras{tenis} \textsuperscript{(34)}.

%%K t02-29-1 p
29. --- Acest din urmă sport este desfătarea elevilor mari. Și profesorii noștri
se prind bucuros la el. Cere multă îndemânare și sprinteneală, și îmi place foarte
mult. \VUgras{Leagănul} \textsuperscript{(35)} îl las fetelor.

%%K t02-30-1 p
30. --- Nu îmi este pe plac nici un alt joc englezesc, care poate ieși
primejdios.

%%K t02-30-1b p
Vorbesc despre \VUgras{fotbal} \textsuperscript{(36)}, bucuria fratelui meu mai
mare. Prefer vechiul nostru joc de-a \VUgras{prinșii} \textsuperscript{(38)}, tot
atât de sănătos și mult mai puțin grosolan.

%%K t02-tit-7 sub
\VUsaut{4.6mm}
\VUpk{10.2pt}{40}{Bicicleta și jocurile cu mingea.}
\VUsaut{3.0mm}

%%K t02-31-1 p
31. --- Și eu ocolesc bucuros pajiștea cu \VUgras{bicicleta}
\textsuperscript{(37)}.

%%K t02-32-1 p
32. --- La anul voi învăța să \VUgras{călăresc} \textsuperscript{(39)}. Ce frumos
este calul când merge frumos la pas sau la trap și trece piedicile în galop!

%%K t02-33-1 p
33. --- Vara învățăm să \VUgras{înotăm} \textsuperscript{(40)}. Ne cufundăm și ne
scăldăm cu desfătare în apa curată și răcoroasă a râului. Uneori la sfârșit
slobozim un \VUgras{balon} de hârtie \textsuperscript{(41)}, care se ridică
măreț în văzduh de îndată ce se umple cu aer cald.

%%K t02-34-1 p
34. --- Mai sunt și multe alte jocuri, dar nu aș sfârși niciodată să le înșir, și
din această scurtă povestire se vede că joile pe frumoasa noastră pajiște nu sunt
deloc plicticoase.

%%K t02-34-2 p
N.~B. --- \textit{Profesorului nu îi va fi greu să întregească acest capitol,
descriind mai amănunțit fiecare dintre jocurile cele mai însemnate.}
\VUsaut{6.0mm}
\VUfilet{22mm}
